%!TEX root = ../hutbthesis_main.tex

\chapter{系统核心实现技术}

\section{整体架构与运行机制}

针对现有CARLA仿真平台API接口分散、功能模块紧耦合、交互链路不统一等问题，本文采用分层解耦的分布式架构思想，构建``用户交互层---MCP服务层---仿真执行层''三级体系结构，以实现功能独立、模块解耦、接口统一、易于扩展的工程目标。

\begin{figure}[H]
\centering
\includegraphics[width=0.95\linewidth,keepaspectratio]{images/image8.png}
\caption{系统整体架构图}
\label{fig:system-architecture}
\end{figure}

如图~\ref{fig:system-architecture}所示，系统整体运行遵循闭环控制机制。用户经由交互层输入自然语言指令或结构化命令，该指令在MCP服务层完成语义解析、路由调度、参数校验与上下文匹配后，被分发至相应的仿真执行单元；仿真器层完成对车辆、无人机、场景、环境等对象的操作后，将执行状态、传感器数据及异常信息沿原路径反馈至交互层进行可视化呈现。该机制确保了指令可控、状态可查、结果可溯，满足自动驾驶仿真对高可靠性与实时性的要求。

用户交互层承担多模态指令输入与结果反馈职能，支持Web界面、语音通道及AI助手三种交互方式，通过统一的消息格式将用户请求向下传递。MCP服务层作为系统的核心调度中枢，集成FastMCP通信框架、工具注册中心、路由分发器、上下文管理器及会话缓存等组件，实现指令的集中式解析与分布式执行。仿真执行层由HUTB/CARLA仿真器与Unreal
Engine编辑器构成，分别负责物理级仿真运算与场景级编辑操作，二者通过独立的MCP服务实例与上层解耦通信，避免资源竞争与单点故障。

\section{FastMCP通信框架与双服务实现技术}

1.双服务并行与端口隔离机制

CARLA仿真控制与Unreal
Engine编辑器操作在资源占用、执行时序及异常域上存在显著差异。若采用单一服务进程同时承载仿真运行时控制与引擎级编辑任务，易因编辑器UI阻塞或蓝图编译延迟导致仿真控制指令丢包或超时。系统为避免仿真控制与引擎编辑发生资源冲突问题，因此使用了HUTB
MCP Server和Unreal MCP
Server两个服务分别部署，这两个服务是互不影响的，独立启动，独立运行，同时都有自己的日志打印功能，其中HUTB
MCP服务监听在2000端口上，用于人在环下的人车仿真进行相关操作。Unreal
MCP服务占据55557端口，主要负责Unreal编辑器相关程序化操作请求，在此基础之上实现仿真与编辑并行操作，互不干扰，提升用户体验。

2.标准化工具封装与注册机制

FastMCP抱着面向工具（Tool -
Oriented）的封装理念，把人车仿真以及虚幻引擎里的全部可控功能抽象成标准工具，各个工具具备专属识别名，结构化的输入参数，规范化的返回值以及统一的异常处理格式，在服务开启的时候会自动注册并动态加载，而且支持运行期间的扩展和热插拔，不需要重新启动主服务就能增添功能模块，这样一种设计使得大模型可以以统一的形式调用车辆，无人机，天气，传感器，Actor，蓝图，UI等各种能力，极大地简化了系统的整合复杂程度。

\begin{figure}[H]
\centering
\includegraphics[width=12.9cm,keepaspectratio]{images/image9.png}
\caption{FastMCP双服务通信机制图}
\label{fig:fastmcp-dual-service}
\end{figure}

在图~\ref{fig:fastmcp-dual-service}中，在用户通过指令经解析调度后，可以分别驱动CARLA自动驾驶仿真与Unreal引擎场景编辑，通过统一数据采集与Web可视化界面，形成协同控制闭环。HUTB
MCP
Server监听2000端口，专责CARLA仿真器的运行时控制，包括车辆/无人机生成、动力学调节、传感器数据采集、天气环境调控及交通流管理；Unreal
MCP Server监听55557端口，专责Unreal
Engine编辑器的程序化操作，包括Actor管理、蓝图编辑、节点图构建、UMG
UI设计及关卡配置。两服务独立启动、独立运行、独立记录日志，互不阻塞，从而在保证仿真控制实时性的同时，支持引擎编辑的复杂长时任务。

3.长连接通信与上下文状态同步。

系统依靠TCP长连接达成低延迟，高稳定的指令通信，其包含心跳保活，断线重连，超时重试，数据校验等可靠性机制，MCP服务层保留全局会话上下文，随时记录并更新仿真环境状态，涉及当前车辆位置，姿态，速度，无人机目标点，天气参数，Actor列表，光照信息，场景对象等内容，上下文不断同步，使得大模型能按照当前真实的仿真环境产生连续，合理又可行的指令序列，免除指令冲突，参数越界以及状态不一致的情况

\section{HUTB/CARLA仿真控制实现技术}

\subsection{仿真世界管理与Actor生命周期控制}

CARLA以World对象为全局调度单元，统一管理场景内所有静态实体与动态Actor的资源分配、状态更新及垃圾回收。本文基于CARLA
Python
API设计Actor全生命周期状态机，涵盖蓝图筛选、条件生成、坐标初始化、姿态设定、属性查询、状态监测、批量销毁及内存释放等阶段。

HUTB/CARLA把World当作核心运作单元，针对场景里所有的静态物和动态体执行统一调度，资源分配以及状态保存，仿真世界覆盖地图环境，光照系统，天气系统，交通管理模块以及全部智能体对象，形成起一个完善的虚拟交通环境。

系统支持按蓝图名称精确生成、按预设坐标定点生成、按随机区域分布生成及批量队列生成四种模式，并提供基于类型标签的分类管理与基于唯一ID的精确定位能力。通过显式维护Actor实例表与引用计数，在销毁阶段强制解除Python端与CARLA服务器端的句柄绑定，防止内存泄漏。该机制可快速构建大规模、多对象、高复杂度的仿真测试场景，为无人车环境感知、行为决策及运动控制算法提供稳定、可复现的仿真运行环境。

在系统当中，车辆，行人，障碍物，交通标志，路灯这些实体都是以Actor对象的形式存在的，每一个Actor都具备独特的编号，空间坐标，旋转角度，缩放比例，物理特性，碰撞体以及生命时限，可以经由API实施单独的操控，并执行批量的运作。

本系统依靠CARLA Python
API达到对Actor全生命时段的控制，遵照蓝图名称实施筛选生成。按照预设坐标执行生成。依照随机区域展开生成。成批执行生成。随时设置姿态。查询属性。监测状态。动态执行销毁并释放内存等职能，该系统具备对Actor实施分类经营，条件甄别，按类型汇总以及按ID准确定位的能力，可以快速创建起大规模，多对象，高复杂度的仿真检测场景，给无人车的环境感知，行为决策以及运动控制算法赋予稳定，可重现，可拓展的仿真运行环境。

\subsection{车辆动力学控制与自动驾驶实现}

车辆动力学模块依托PhysX物理引擎构建，需模拟油门开度与加速度的非线性映射、转向比与转向迟滞、制动压力与减速度关系、悬挂行程与车身姿态、轮胎滑移率与路面摩擦系数等关键动力学特性。本文对车辆控制接口进行高度封装，对外提供车辆创建、坐标初始化、姿态设定、手动操控、自动驾驶模式切换、速度限制、灯光控制及销毁等原子操作接口。

在自动驾驶模式下，车辆由CARLA内置的Traffic
Manager统一调度，通过全局路径规划、车道保持、动态避障、红绿灯识别、跟车行驶及变道转弯等模块实现城市道路行驶逻辑的自动复现。该模式可直接用于无人车决策规划算法的验证与对比评估。通过将PhysX动力学仿真与Traffic
Manager自动驾驶逻辑深度融合，系统能够稳定复现各类常见行驶工况与危险场景，提升算法验证的安全性与覆盖度。

\subsection{多类型传感器数据采集与输出}

传感器系统对于无人车的感知能力十分关键，该系统可在车辆或者无人机上动态地挂载，设置并启动或者关闭各类虚拟传感器，涉及RGB相机，深度相机，激光雷达，毫米波雷达，IMU，GNSS碰撞检测器以及车道线检测器等。

各类传感器均能自定义分辨率，视场角，探测距离，刷新频率，安装位置以及安装角度，其数据输出格式同真实无人车硬件相契合，可输出图像数据，三维点云，目标检测框，位置坐标，角速度，加速度，航向角等信息。

传感器数据会按照预定频率即时反馈到MCP服务层，这些数据可被用于感知算法的验证，数据集的形成，目标检测模型的考察，环境认知以及决策系统的输入部分，经由对统一接口实施封装处理，用户不必理会底层驱动存在的差别，就能很快达成多传感器的设置并获取相关数据，大幅缩减创建仿真环境及其调试过程中的成本。

\subsection{天气与环境参数化调控}

天气与环境系统具备多参数连续可调的功能，可以精确模拟各类气象，光照以及时间因素给无人车感知和控制带来的影响，该系统允许针对云量，雨量，雾量，积水深度，湿度，灰尘浓度，光照强度，太阳高度角，时间段（早晚夜，黄昏）等参数单独或者综合执行设定。

系统内置多种常见天气模式以便操作方便，晴天、阴天、小雨、中雨、大雨、暴雨、大雾、傍晚、夜晚、沙尘暴等都包括，在此基础上可以自定义特殊天气组合来模拟如强降雨、低能见度、强逆光、湿滑路面等情况，以全面测试无人驾驶车辆对于各种不良条件下的鲁棒性和安全性。

首先，建立自然语言气象语义到参数集合的映射规则库。例如，``暴雨''对应\{precipitation:
100, fog\_density: 50, wind\_intensity: 30, cloudiness:
90\}，``黄昏''对应\{sun\_altitude\_angle: 10, sun\_azimuth\_angle: 270,
cloudiness:
20\}。该映射规则库支持内置常见模式（晴天、阴天、小雨、大雾、夜晚、沙尘暴等）与用户自定义组合。

其次，设计基于线性插值的天气渐变过渡算法，避免环境参数瞬时跳变导致的仿真不连续。设目标天气参数为W\_target，当前参数为W\_current，过渡时长为T，刷新周期为Δt，则第k步的中间状态为：

\begin{equation}
W_k=W_{current}+(W_{target}-W_{current})\times\frac{k\Delta t}{T}
\label{eq:weather-transition}
\end{equation}

其中k = 1, 2, ..., T/Δt。本文取Δt = 0.1 s，默认过渡时长T = 2.0
s，实现视觉与物理层面的平滑切换。

此外，系统建立环境参数与车辆物理特性的联动机制。例如，路面摩擦系数$\mu$随降水量$p$增加而衰减，其近似关系如式~\eqref{eq:rain-friction}所示：
\begin{equation}
\mu=\mu_0(1-0.003p)
\label{eq:rain-friction}
\end{equation}
其中$\mu_0$为干燥路面基准摩擦系数。该机制使得天气调控不仅影响视觉渲染，更直接作用于车辆动力学响应，提升了仿真环境的物理真实性。

系统经由对天气和环境实施参数化控制，可以形成多种化，高覆盖度的仿真场景库，从而满足无人车在不同地域，不同季节，不同气象条件下的算法验证需求，改良仿真检测的真实性和全面性。

\section{自然语言指令解析与映射实现技术}

1.意图结构化解析与要素提取

自然语言认识模块把用户所输入的模糊文本指令分解成操作对象，动作类型以及参数列表这三大核心要素，操作对象是车辆，动作类型包含生成和自动驾驶两个部分，其中生成对应的是创建新对象的动作，而自动驾驶则属于一种功能设置或者说是状态改变，参数方面有车型即轿车，坐标点位于（100,200,0），还有自动驾驶是否被激活这个布尔值参数，经过结构化解析之后，那些含糊不清的语言就变成了计算机能够执行下去的具体逻辑规则。

2.模糊参数标准化映射

系统建立语义-数值映射库以解决自然语言所具有的不确定性，该库可以将自然语言中的描述性词语转化为仿真能够理解的标准参数，``广场''对应预设坐标，``暴雨''则为雨量=100、能见度=50、风速=30，``自动驾驶''即使能开关=TRUE，``白天''亦为时间参数14:00等，这些参数对应关系可以使命令被直接执行。

3.意图---工具---API三级路由

系统创建起意图-工具-
API这三级路由表，从而达成从自然语言到底层调用的完整映射效果，经过分析得到的意图会自动与相应的MCP工具相契合，该工具接着去调用对应仿真或者编辑器的API来完成指令的执行，这个路由机制包含车辆，无人机，天气，传感器，场景，Actor，蓝图，UI等所有功能领域，做到了全场景下由自然语言来激发控制。

\section{模块化工具集设计与实现技术}

1.工具集统一设计规范

所有的功能工具都应遵循一套统一的标准，有唯一的名字，参数明确，格式一致，结果一样，异常可被捕获，日志可查询，这样可以使程序具有高内聚，低耦合性，便于后续维护、扩展及合作开发。2.HUTB仿真控制工具集

HUTB工具集包括六种类型，即车辆控制工具、无人机飞行工具、天气环境工具、传感器管理工具、场景编辑工具和交通参与者工具，每一种工具都提供了细粒度操作接口，在整个仿真的过程中都可以进行相关操作。

3.Unreal编辑器控制工具集

Unreal工具集包含六类工具，分别是Actor管理工具，蓝图编辑工具，节点图编辑工具，UMG
UI工具，关卡设置工具以及项目设置工具，它支持编辑器全功能的远程调用。

4.统一日志与异常处理机制

系统达成全局统一的日志输出，错误码定义，异常捕捉，指令重试，状态回滚机制，各个工具若执行失败，会返回标准化的错误信息，这有益于调试和问题定位，可以加强系统的整体健壮性。

\enlargethispage{2\baselineskip}
\section{轻量化部署与自动化启动技术}

1.Python虚拟环境隔离

系统利用Python
venv虚拟环境来达成依赖隔离，防止不同项目以及不同库版本相互冲突，这个环境具备可打包性，可迁移性，而且可以一键复原，极大地优化了部署效率。

2.自动化部署与一键启动

本文设计基于依赖拓扑排序的多进程自动编排算法，实现系统的一键化部署。部署脚本（start.bat）依次执行以下操作：检测Python环境→创建或激活虚拟环境→安装缺失依赖→检查CARLA服务端可执行文件→按用户选择启动CARLA服务器（支持DirectX11、最小配置、无头模式或跳过仿真仅启用AI对话）→启动HUTB
MCP Server→启动Unreal MCP Server→启动Web可视化服务→自动打开浏览器。

脚本内置错误检测与状态判断逻辑：若某一步骤失败（如端口占用、依赖安装超时、CARLA进程崩溃），立即输出标准化错误提示并终止后续步骤，防止级联故障。该方案将传统手动部署的十余个步骤压缩为单次点击，显著降低了系统的使用门槛与部署周期。代码~\ref{lst:start-script}给出了启动脚本的流程化表达。

\begin{lstlisting}[style=hutbcode,language=bash,basicstyle=\ttfamily\tiny,numberstyle=\tiny\color{gray},aboveskip=1pt,belowskip=1pt,framesep=1pt,lineskip=-1pt,caption={一键启动脚本流程示例},label={lst:start-script}]
@echo off
call .venv\Scripts\activate
python -m pip install -r requirements.txt
start carla_server.exe -dx11
start python hutb_mcp_server.py
start python unreal_mcp_server.py
start python web_app.py
\end{lstlisting}

3.外部配置文件管理

系统将服务端口号、CARLA安装路径、日志级别、仿真默认参数、启动选项等纳入外部配置文件（config.yaml）进行集中管理。用户无需修改源码即可调整系统行为，提升了系统的灵活性、可配置性与跨平台适配能力。配置文件采用分层结构，区分开发环境、测试环境与生产环境，支持多场景快速切换。配置文件的基本组织方式如代码~\ref{lst:config-yaml}所示。

\begin{lstlisting}[style=hutbcode,language=bash,basicstyle=\ttfamily\tiny,numberstyle=\tiny\color{gray},aboveskip=1pt,belowskip=1pt,framesep=1pt,lineskip=-1pt,caption={config.yaml配置结构示例},label={lst:config-yaml}]
hutb_mcp_port: 2000
unreal_mcp_port: 55557
carla_path: D:/CARLA
log_level: INFO
web:
  host: 127.0.0.1
  port: 8000
\end{lstlisting}
